\chapter{Introduction}

[Need to rewrite as writeup is now quite different]

The problem of correctly assigning the price of an option on a stock is a widely-researched field and is essential in minimising market inefficiencies. Black and Scholes have developed a closed-form solution to price the European put and call options \cite{black1973pricing}, but their theory can be extended to American options as well. American options allow holders of the contract to exercise the option at any time before expiry, which turns our pricing problem into an optimal stopping problem. 

In this study, we will specifically focus on put options on stocks paying no dividends, as it's well-known that it is always optimal to wait until expiry to exercise a call option provided there are no dividends [need to cite], and hence the American and European prices coincide. For an asset $S$, the put option has the payoff at expiry
\[g(S) = (K-S)^+\]
where $K$ is the pre-determined strike price.

[Talk about the structure of the paper here (to be completed after more writeup is done)]

\section{Probabilistic Setup}

\subsection{Risk-neutral measure and arbitrage-free pricing}

The market in the real world is not arbitrage free, however, so there are implicit assumptions that have been made that we assume hold across the models that we'll explore:
\begin{itemize}
    \item Shares of stock can be infinitely subdivided
    \item Interest rates for investing and borrowing are the same (in reality, the rate for borrowing is generally higher)
    \item The selling and buying prices coincide (i.e. there is no bid-ask spread)
    \item Buying and selling causes no friction in the market (no transaction costs, taxes, etc)
    \item A trader can long and short any amount of one stock, and borrow any amount from the bank
\end{itemize}
